\documentclass[acmsmall,screen,review]{acmart}

\usepackage{proof}
\usepackage{fancyvrb}
\usepackage{minted}
\setminted{fontsize=\small} % Remove this for larger code snippets

\usepackage{tikz}
\usepackage{graphicx}
\usepackage{mathpartir}
\usepackage{color}
\usepackage{pgflibraryarrows}
\usepackage{pgflibraryshapes}
\usepackage{pgfbaseimage}
\usepackage{pgfplots}

\usetikzlibrary{decorations.text}
\usetikzlibrary{trees}
\usetikzlibrary{fadings}

\usepackage{xcolor}

\usepackage{amsthm}

\let\Bbbk\relax
\usepackage{amssymb}

\let\Bbbk\relax
\usepackage{amssymb}

\usepackage[ruled,vlined]{algorithm2e}
\SetAlgoLined
\SetAlgoNoEnd
\SetArgSty{textup}
\SetKwProg{Program}{Program}{\hspace{0.3em}:}{}
\SetKwInOut{Input}{input}\SetKwInOut{Output}{output}
\SetKwIF{Match}{ElseMatch}{Otherwise}{match}{with}{else}{else}{endmatch}

\usepackage{fancyhdr}

\usepackage{booktabs}

\usepackage{multicol}

\setcopyright{acmlicensed}
\copyrightyear{2026}
\acmYear{2026}
\acmDOI{XXXXXXX.XXXXXXX}

\acmJournal{JACM}
\acmVolume{37}
\acmNumber{4}
\acmArticle{111}
\acmMonth{8}

%%
%% Submission ID.
%% Use this when submitting an article to a sponsored event. You'll
%% receive a unique submission ID from the organizers
%% of the event, and this ID should be used as the parameter to this command.
%%\acmSubmissionID{123-A56-BU3}


%%
%% The majority of ACM publications use numbered citations and
%% references.  The command \citestyle{authoryear} switches to the
%% "author year" style.
%%
%% If you are preparing content for an event
%% sponsored by ACM SIGGRAPH, you must use the "author year" style of
%% citations and references.
%% Uncommenting
%% the next command will enable that style.
%%\citestyle{acmauthoryear}


\newcommand{\abs}[4]{{#1}\, #2\! : \! #3.\, #4}
\newcommand{\absu}[3]{{#1}\, #2.\, #3}
\mathchardef\mhyph="2D % Define a "math hyphen"

% unicode
\usepackage[utf8]{inputenc}

\DeclareUnicodeCharacter{2261}{\ensuremath{\equiv}}
\DeclareUnicodeCharacter{2293}{\ensuremath{\sqcap}}
\DeclareUnicodeCharacter{21D3}{\ensuremath{\Downarrow}}
\DeclareUnicodeCharacter{226A}{\ensuremath{\ll}}
\DeclareUnicodeCharacter{2200}{\ensuremath{\forall}}
\DeclareUnicodeCharacter{2124}{\ensuremath{\mathbb{Z}}}
\DeclareUnicodeCharacter{1D539}{\ensuremath{\mathbb{B}}}
\DeclareUnicodeCharacter{2115}{\ensuremath{\mathbb{N}}}
\DeclareUnicodeCharacter{2A7D}{\ensuremath{\le}}
\DeclareUnicodeCharacter{3BB}{\ensuremath{\lambda}}
\DeclareUnicodeCharacter{3B4}{\ensuremath{\delta}}
\DeclareUnicodeCharacter{3C9}{\ensuremath{\omega}}
\DeclareUnicodeCharacter{3B3}{\ensuremath{\gamma}}
\DeclareUnicodeCharacter{3C4}{\ensuremath{\tau}}
\DeclareUnicodeCharacter{21D2}{\ensuremath{\To}}
\DeclareUnicodeCharacter{22B2}{\ensuremath{\lhd}}
\DeclareUnicodeCharacter{2080}{\ensuremath{_0}}
%\definecolor{darkred}{RGB}{100,10,10}


\newcommand{\BD}[1]{\textcolor{olive}{[BD: #1]}}
\newcommand{\PA}[1]{\textcolor{cyan}{[PA: #1]}}

\newcommand{\utp}[0]{\mathcal{U}}
\newcommand{\lin}[0]{\lambda^{\iota\nu}}
\newcommand{\interp}[1]{\llbracket #1 \rrbracket}
\newcommand{\interpl}[0]{\llbracket}
\newcommand{\interpr}[0]{\rrbracket}
\newcommand{\choice}[0]{\xi}
\newcommand{\simto}[0]{\triangleright}
\newcommand{\elcap}[0]{\cap}

\newcommand{\dedu}[0]{\triangleright}
\newcommand{\ctxtup}[1]{\ulcorner #1 \urcorner}
\newcommand{\ctxtdn}[1]{\llcorner #1 \lrcorner}
\newcommand{\ctxtmv}[1]{\langle #1 \rangle}

\newcommand{\case}[1]{\noindent\underline{Case #1:}}
\newcommand{\caseb}[0]{\noindent\underline{Case:}}

\newcommand{\pcase}[1]{\noindent\underline{#1:}}


\newcommand{\adbmal}[0]{\reflectbox{$\lambda$}}

\newcommand{\casetri}[0]{\RHD}
\newcommand{\lke}[0]{\lhd}

\newcommand{\lam}[2]{\lambda\, #1.\, #2}
\newcommand{\mal}[2]{\abdmal\, #1.\, #2}
\newcommand{\all}[2]{\forall\, #1.\, #2}
\newcommand{\Lam}[2]{\Lambda\, #1.\, #2}
\newcommand{\thereis}[2]{\exists\, #1.\, #2}
\newcommand{\nuu}[2]{\nu\, #1.\,#2}
\newcommand{\inn}[2]{\langle #1 \rangle_{#2}}
\newcommand{\exel}[2]{\adbmal\,#1.\,#2}
\newcommand{\recc}[2]{\textit{rec}\, #1.\, #2}
\newcommand{\mew}[2]{\mu_{#1}\, #2}
\newcommand{\convl}[0]{\blacktriangleleft}
\newcommand{\convr}[0]{\blacktriangleright}
\newcommand{\conjin}[0]{\smalltriangleleft}
\newcommand{\conjout}[0]{\smalltriangleright}
\newcommand{\dcdot}[0]{\mathbin{\ast}}
\newcommand{\To}[0]{\Rightarrow}
\newcommand{\releq}[0]{\topdoteq}

\newcommand{\omga}[4]{\omega\ #1(#2)\,:\,#3.\ #4}
\newcommand{\clause}[2]{#1 \to #2}

\newcommand{\hysep}[0]{\mathbin{\blacklozenge}}

\newcommand{\todo}[1]{\textbf{[#1]}}
\newcommand{\rref}[1]{\{\ref{#1}\}}

\newtheorem{theorem}{Theorem}
\newtheorem{lemma}[theorem]{Lemma}
\newtheorem{corollary}[theorem]{Corollary}
\newtheorem{definition}[theorem]{Definition}
\newtheorem{notation}[theorem]{Notation}
\newtheorem{proposition}[theorem]{Proposition}


%\newtheorem{definition}[theorem]{Definition}


% \newenvironment{mycomment}
%     {\color{blue} \it
%     }
%     {
%     }
\newcommand{\shift}[2]{{\uparrow^{#1}_{#2}}}
\newcommand{\textblue}[1]{\textcolor{blue}{#1}}


\newcommand{\llpic}[0]{\begin{tikzpicture}
  \path (1,1) node(a){DATA tt INT};
  \path (4,1) node(b){DATA tt BOOL};
  \path (7,1) node(c){FAIL tt};
  \path (1,3) node(aa){DATA ff INT};
  \path (4,3) node(bb){DATA ff BOOL};
  \path (9,3) node(cc){FAIL ff};
  \path (4,-1) node(u){UNKNOWN};
  \path (7,3) node(l){LOOP};

  \draw[arrows = {-Stealth[length=5pt]}] (a.north) -- (aa.south);
  \draw[arrows = {-Stealth[length=5pt]}] (b.north) -- (bb.south);
  \draw[arrows = {-Stealth[length=5pt]}] (c.north) -- (cc.south);    
  \draw[arrows = {-Stealth[length=5pt]}] (u.north) -- (a.south);
  \draw[arrows = {-Stealth[length=5pt]}] (u.north) -- (b.south);
  \draw[arrows = {-Stealth[length=5pt]}] (u.north) -- (c.south);    
  \draw[arrows = {-Stealth[length=5pt]}] (a.north) -- (l.south);
  \draw[arrows = {-Stealth[length=5pt]}] (b.north) -- (l.south);
  \draw[arrows = {-Stealth[length=5pt]}] (c.north) -- (l.south);    
  
\end{tikzpicture}}


%%
%% end of the preamble, start of the body of the document source.
\begin{document}

\title{There Must Be Some Kind of Mistake: Distinguishing Errors from Uncertainty in Type Checking}


%% \author{Aaron Stump}
%% \email{stumpaa@bc.edu}
%% \affiliation{%
%%   \institution{Boston College}
%%   \city{Boston}
%%   \state{Massachusetts}
%%   \country{USA}
%% }


%%
%% By default, the full list of authors will be used in the page
%% headers. Often, this list is too long, and will overlap
%% other information printed in the page headers. This command allows
%% the author to define a more concise list
%% of authors' names for this purpose.
\renewcommand{\shortauthors}{\ }%Blanchette et al.}

%%
%% The abstract is a short summary of the work to be presented in the
%% article.
\begin{abstract}
  This paper extends work of Garby, Bahr, and Hutton on calculating
  type checkers, to distinguish errors that will certainly occur at
  runtime from uncertainty due to the approximate nature of type
  checking.  The work also extends this to a simple situation with
  possibly diverging programs, where again we distinguish certainty of
  divergence from uncertainty due to approximation.
\end{abstract}

%%
%% The code below is generated by the tool at http://dl.acm.org/ccs.cfm.
%% Please copy and paste the code instead of the example below.
%%
\begin{CCSXML}
<ccs2012>
   <concept>
       <concept_id>10011007.10011074.10011099.10011692</concept_id>
       <concept_desc>Software and its engineering~Formal software verification</concept_desc>
       <concept_significance>500</concept_significance>
       </concept>
   <concept>
       <concept_id>10003752.10010124.10010138.10010142</concept_id>
       <concept_desc>Theory of computation~Program verification</concept_desc>
       <concept_significance>300</concept_significance>
       </concept>
   <concept>
       <concept_id>10003752.10003790.10011740</concept_id>
       <concept_desc>Theory of computation~Type theory</concept_desc>
       <concept_significance>500</concept_significance>
       </concept>
 </ccs2012>
\end{CCSXML}

\ccsdesc[500]{Software and its engineering~Formal software verification}
\ccsdesc[300]{Theory of computation~Program verification}
\ccsdesc[500]{Theory of computation~Type theory}

\keywords{type checking, nontermination, type errors}

%\received{20 February 2007}
%\received[revised]{12 March 2009}
%\received[accepted]{5 June 2009}

\theoremstyle{definition}
\newtheorem{specification}[section]{Specification}

%%
%% This command processes the author and affiliation and title
%% information and builds the first part of the formatted document.
\maketitle

\section{Introduction}

\citet{garby+25} presents the derivation of a type checker for a
simple expression language, in the style of program calculation.  One
essentially tries to prove a specification about code that has yet to
be determined, and discovers the definitions (or constraints on them)
necessary in the course of attempting to complete the proof.  Their
specification states that the type computed by the type checker
\verb|texp| for an \textit{Expr} $e$ is an approximation to the true type
of $e$.  The true type is the type computed (by a function
\textit{tval}) for the value of $e$, as produced by an evaluator
\textit{eval}.  The notion of approximation is that the \textit{ERROR}
type returned by the type checker to signal a type error approximates
any of the types for values.  With the types shown in Figure~\ref{fig:orig},
their specification is
\[
\textit{texp}\ e \le \textit{tval}\ (\textit{eval}\ e)
\]

\begin{figure}
\begin{minted}{haskell}
data Type = INT | BOOL | ERROR
\end{minted}
\caption{Definitions from \citet{garby+25}}xo
\label{fig:orig}
\end{figure}





\bibliographystyle{ACM-Reference-Format}
\bibliography{paper.bib}


\end{document}
% \endinput
%%
%% End of file `sample-acmsmall-submission.tex'.
